\usepackage{graphicx, amsmath, amssymb, amsthm, amsfonts}
\usepackage[mathscr]{eucal}
\usepackage{thmtools} % without this cref says theorem for everything
\usepackage{microtype}
\usepackage{tikz-cd}
\usepackage{xparse}         % better newcommand - NewDocumentCommand
\usepackage{verbatim} % comment environment, handy for debugging errors
\usepackage{xy}
\usepackage[all]{xypic}
\usepackage{spectralsequences}
\usepackage{mathtools}
\usepackage[textsize=tiny]{todonotes}
% keep these two last to avoid weird errors
\usepackage[pagebackref, colorlinks, citecolor=Red, linkcolor=NavyBlue, urlcolor=NavyBlue]{hyperref}
\usepackage[capitalise, nameinlink]{cleveref}
\hypersetup{linktoc=all} %set to all if you want both sections and subsections linked

% Noah's notes
\newcommand{\noahnote}[1]{\todo[color=blue!40,linecolor=blue!40!black,size=\tiny]{Noah: #1}}
\newcommand{\noahnoteil}[1]{\ \todo[inline,color=blue!40,linecolor=blue!40!black,size=\normalsize]{Noah: #1}}

% David's notes
\newcommand{\dcnote}[1]{\todo[color=yellow!40,linecolor=yellow!40!black,size=\tiny]{David: #1}}
\newcommand{\davidnote}[1]{\todo[color=yellow!40,linecolor=yellow!40!black,size=\tiny]{David: #1}}
\newcommand{\davidnoteil}[1]{\ \todo[inline,color=yellow!40,linecolor=yellow!40!black,size=\normalsize]{David: #1}}

% Inbar's notes
\newcommand{\inbarnote}[1]{\todo[color=green!40,linecolor=green!40!black,size=\tiny]{Inbar: #1}}
\newcommand{\inbarnoteil}[1]{\ \todo[inline,color=green!40,linecolor=green!40!black,size=\normalsize]{Inbar: #1}}

% Marc's notes
\newcommand{\marcnote}[1]{\todo[color=red!40,linecolor=red!40!black,size=\tiny]{Marc: #1}}
\newcommand{\marcnoteil}[1]{\ \todo[inline,color=red!40,linecolor=red!40!black,size=\normalsize]{Marc: #1}}


\author[Chan]{David Chan}
\address[Chan]{Department of Mathematics, Michigan State University, East Lansing, MI 48824 }
\email{chandav2@msu.edu}

\author[Gotliboym]{Marc Gotliboym}
\address[Gotliboym]{Department of Mathematics, Michigan State University, East Lansing, MI 48824 }
\email{marcg57@msu.edu}

\author[Klang]{Inbar Klang}
\address[Klang]{Department of Mathematics,
Vrije Universiteit Amsterdam - Faculty of Science,
De Boelelaan 1111,
1081 HV Amsterdam,
The Netherlands}
\email{i.klang@vu.nl}

\author[Wisdom]{Noah Wisdom}
\address[Wisdom]{Department of Mathematics,
         Northwestern University,
         Evanston, IL, 60208
         USA}
\email{NoahAnkney2026@u.northwestern.edu}


\newcommand{\C}{\mathbb{C}}
\newcommand{\R}{\mathbb{R}}
\newcommand{\Z}{\mathbb{Z}}
\newcommand{\A}{\mathcal{A}}
\newcommand{\F}{\mathbb{F}}
\newcommand{\E}{\mathbb{E}}
\newcommand{\N}{\mathbb{N}}
\DeclareMathOperator{\Hom}{Hom}
\DeclareMathOperator{\colim}{colim}
\DeclareMathOperator{\id}{id}
\DeclareMathOperator{\Sym}{Sym}

\newcommand{\Tamb}{\mathsf{Tamb}}
\newcommand{\Mack}{\mathsf{Mack}}
\newcommand{\CRing}{\mathsf{CRing}}
\newcommand{\Alg}{\mathsf{Alg}}
\newcommand{\FP}{\mathrm{FP}}
\newcommand{\ETHH}{\mathrm{ETHH}}
\newcommand{\ethh}{\Res^{C_p\times S^1}_{C_p\times 1}\mathrm{ETHH}}
\newcommand{\THH}{\mathrm{THH}}
\newcommand{\HH}{\mathrm{HH}}
\newcommand{\ETC}{\mathrm{ETC}}
\newcommand{\TC}{\mathrm{TC}}
\newcommand{\cycsp}{\mathrm{CycSp}}
\newcommand{\Sp}{\mathrm{Sp}}


%\theoremstyle{definition}
\newtheorem{theorem}{Theorem}[section]
\newtheorem{lemma}[theorem]{Lemma}

\newtheorem{proposition}[theorem]{Proposition}
\newtheorem{corollary}[theorem]{Corollary}
\newtheorem{conjecture}[theorem]{Conjecture}

\newtheoremstyle{BoldRemark} % name
{10pt}                    % Space above
{10pt}                    % Space below
{\upshape}                   % Body font
{}                           % Indent amount
{\bfseries}                  % Theorem head font
{.}                          % Punctuation after theorem head
{.5em}                       % Space after theorem head
{}  % Theorem head spec (can be left empty, meaning ‘normal’)
\theoremstyle{BoldRemark}
\newtheorem{remark}[theorem]{Remark}
\newtheorem{definition}[theorem]{Definition}
\newtheorem{example}[theorem]{Example}
\newtheorem{notation}[theorem]{Notation}

\newenvironment{cd}{
  \begin{center}\begin{tikzcd}}{
    \end{tikzcd}\end{center}}

\newcommand{\<}{\left\langle}
\renewcommand{\>}{\right\rangle}


\newcommand{\mac}[1]{\underline{#1}} % 'mackey' ie underline everything
\newcommand{\dual}{\ensuremath{^{\mathrm{op}}}}
\newcommand{\ginfcat}{\ensuremath{G}-\ensuremath{\infty}-category}
\DeclareRobustCommand{\infcat}{%
  \ifmmode
   \mathrm{Cat}_\infty%
  \else
   \ensuremath{\infty}-category%
  \fi
}

\newcommand{\inv}{^{-1}}
\newcommand{\cofib}{\operatorname{cofib}}
\newcommand{\fib}{\operatorname{fib}}
\newcommand{\can}{\mathrm{can}} % canonical maps
\newcommand{\ev}{\mathrm{ev}} % evaluation
\newcommand{\we}{\simeq} % weakly equivalent
\newcommand{\aut}{\operatorname{Aut}}
\newcommand{\infloop}{\Omega^\infty}
\newcommand{\infsusp}{\Sigma^\infty_+}
\DeclarePairedDelimiter{\ceil}{\lceil}{\rceil}
\DeclarePairedDelimiter{\floor}{\lfloor}{\rfloor}
\newcommand{\kernel}{\ensuremath{\operatorname{Ker }}}
\newcommand{\smsh}{\wedge}
\newcommand{\cyc}{\mathrm{cyc}}

\newcommand{\Res}{\operatorname{Res}}
\newcommand{\triv}{\mathrm{triv}}
\newcommand{\Pic}{\operatorname{Pic}}



\newcommand{\fun}{\mathrm{Fun}}
\newcommand{\Fun}{\mathrm{Fun}}
\NewDocumentCommand{\Spg}{O{G}}{\Sp_{#1}}
\NewDocumentCommand{\mspg}{O{G}}{\mac{\Sp}_{#1}}
\newcommand{\mfun}{{\mac{\fun}}}
\newcommand{\map}{{\mathrm{map}}}


\NewDocumentCommand{\etc}{ O{G} }{\mac{\mathrm{ETC}}}

% parametrized homotopy orbits/fixed points/tate
\NewDocumentCommand{\horb}{O{G}}{_{h#1}}
\NewDocumentCommand{\hfix}{O{G}}{^{h#1}}
\NewDocumentCommand{\tate}{O{G}}{^{t#1}}
\NewDocumentCommand{\geo}{O{G}}{^{\Phi #1}}
\newcommand{\horbp}{\horb[C_p]}
\newcommand{\hfixp}{\hfix[C_p]}
\newcommand{\tatep}{\tate[C_p]}
\newcommand{\geop}{\geo[C_p]}


\NewDocumentCommand{\thfix}{O{G} O{N}}{^{h_{#1}#2}} %twisted homotopy fix points
\NewDocumentCommand{\thorb}{O{G} O{N}}{_{h_{#1}#2}} %twisted homotopy orbits
\NewDocumentCommand{\ttate}{O{G} O{N}}{^{t_{#1}#2}}
\NewDocumentCommand{\macttate}{O{G} O{N}}{^{\mac{t}_{#1}#2}}
\NewDocumentCommand{\macthfix}{O{G} O{N}}{^{\mac{h}_{#1}#2}}
\NewDocumentCommand{\macthorb}{O{G} O{N}}{_{\mac{h}_{#1}#2}}
\NewDocumentCommand{\macthsp}{O{G} O{N}}{{\mac\Sp}\thfix[#1][#2]}
\NewDocumentCommand{\thsp}{O{G} O{N}}{{\Sp}\thfix[#1][#2]}
\NewDocumentCommand{\hsp}{O{G}}{{\Sp}\hfix[#1]}


\newcommand{\gtatep}{{\ttate[G][C_p]}}
\newcommand{\mgtatep}{{\macttate[G][C_p]}}
\newcommand{\mghfixp}{\macthfix[G][C_p]}
\newcommand{\mghorbp}{\macthorb[G][C_p]}
\newcommand{\mgtatecirc}{{\macttate[G][S^1]}}
\newcommand{\mghfixcirc}{\macthfix[G][S^1]}
\newcommand{\mghorbcirc}{\macthorb[G][S^1]}
\newcommand{\msphgp}{\macthsp[G][C_p]}
\newcommand{\msphgcirc}{\macthsp[G][S^1]}
\newcommand{\msphpcirc}{\macthsp[C_p][S^1]}
\newcommand{\sphpcirc}{\thsp[C_p][S^1]}
\newcommand{\HFp}{H\underline{\mathbb{F}}_p}
\newcommand{\HFnop}{H\underline{\mathbb{F}}}
\newcommand{\Tor}{\mathrm{Tor}}
\newcommand{\Top}{\mathrm{Top}}



% G means optional but delimited by {}
\newcommand{\LEq}[3]{\mathrm{LEq}_{#1:#2}\left(#3\right)}
\NewDocumentCommand{\gLEq}{G{F} G{F'} G{\mathscr{C}} }{%
  \mac{\mathrm{LEq}}_{{#1}:{#2}}\left(#3\right)}


% letter theorems for the intro
\theoremstyle{plain}
\newtheorem{introtheorem}{Theorem}
\renewcommand{\theintrotheorem}{\Alph{introtheorem}} 
\crefname{introtheorem}{Theorem}{Theorems}